%% Le dipôle résistif : convention récepteur, U et I fléchés, la loi d'Ohm en rappel.
\documentclass[tikz,border=4pt]{standalone}
\usepackage{liaisons}
\begin{document}
\begin{tikzpicture}[x=1cm,y=1cm,
  fil/.style={line width=1pt, draw=black!75},
  fleche/.style={-{Stealth[length=5pt,width=4.5pt]}, line width=0.9pt}]

%% le dipôle
\draw[fil, fill=solideB!10] (-0.75,-0.4) rectangle (0.75,0.4);
\node[font=\large] at (0,0) {$R$};
\draw[fil] (-2.1,0) -- (-0.75,0);
\draw[fil] (0.75,0) -- (2.1,0);

%% le courant, dans le même sens des deux côtés
\draw[fleche, draw=solideA] (-1.85,0.28) -- (-1.25,0.28);
\node[font=\small, text=solideA, anchor=south] at (-1.55,0.3) {$I$};
\draw[fleche, draw=solideA] (1.25,0.28) -- (1.85,0.28);
\node[font=\small, text=solideA, anchor=south] at (1.55,0.3) {$I$};

%% la tension, flèche de la borne moins vers la borne plus (convention récepteur)
\draw[fleche, draw=solideE] (1.6,0.95) -- (-1.6,0.95);
\node[font=\small, text=solideE, anchor=south] at (0,0.98) {$U$};

%% ce que l'énoncé ne doit pas donner
\node[anchor=north, font=\small, align=center, inner sep=4pt] at (0,-0.65) {%
  \rappel{$U = R\times I$}\\[3pt]
  \footnotesize $U$ en volts (V), $R$ en ohms ($\Omega$), $I$ en ampères (A)};
\end{tikzpicture}
\end{document}
